\section{Discussion}
This worked example illustrates a direct connection between balanced design,
prespecified quality control, model-based estimates, multiplicity correction,
and sensitivity analysis. Because the ground truth is known, successful recovery
can be assessed separately from biological interpretation. The example also
shows why reporting only a list of significant features is insufficient: effect
sizes, uncertainty, null calibration, and robustness each answer a different
question.

The largest performance loss occurred when a known batch component was omitted.
This observation is not evidence about any biological system; it is a design
property of the simulation. It nevertheless demonstrates a general reporting
principle: technical factors anticipated during design should remain visible in
the analysis and methods. Conversely, the close agreement between two reasonable
normalization strategies indicates that the strongest synthetic signals were
not artifacts of one scaling convention.

The manuscript is intentionally self-contained. A real study should replace the
synthetic resources with persistent repository accessions, report organism- or
cell-line-specific details, identify every exclusion and replicate, and provide
figure-level sample sizes and statistical tests. The Key Resources Table should
list only resources essential to reproducing the reported work, using stable
identifiers wherever possible.
